% !TeX root = 02_intro_ml_polytech.tex
% !TeX encoding = UTF-8
% !TeX spellcheck = en_GB
\begin{frame}[fragile]{A Simple Exercise}{\includegraphics[height=12pt]{01_slides/img/tool.pdf}Classifying some data\dots}
    \begin{enumerate}
        \item \alert<2>{\includegraphics[height=12pt]{01_slides/img/pen-tool.pdf}~Go back to your \texttt{linear\_models.ipynb} notebook}
        \item \alert<3-4>{\includegraphics[height=12pt]{01_slides/img/pen-tool.pdf}~Generate classification data in $K$ different clusters (e.g.\ $K = 4$) in 2D:}
              \begin{itemize}
                  \item Create a vector of $K$ 2D points (they will be the centres of the clusters)
                  \item Create a vector of $250$ labels for each cluster
                  \item Create Gaussian-distributed data as \enquote{center(k) + random noise} for all labels
                  \item Visualise the result using a scatter plot
              \end{itemize}
        \item \alert<5>{\includegraphics[height=12pt]{01_slides/img/pen-tool.pdf}~Split the data in train and test (remember to shuffle!)}
        \item \alert<6>{\includegraphics[height=12pt]{01_slides/img/pen-tool.pdf}~Use, for instance, \texttt{statsmodels}' \textbf{MNLogit} (\emph{multinomial logit}) to fit your data and print the summary}
        \item \alert<7-8>{\includegraphics[height=12pt]{01_slides/img/pen-tool.pdf}~Compute the confusion matrix of the classifier}
        \item \alert<9>{\includegraphics[height=12pt]{01_slides/img/pen-tool.pdf}~Compute accuracy, precision, and recall of the classifier}
              \begin{itemize}
                  \item \includegraphics[height=12pt]{01_slides/img/brain.pdf}~This is a \textbf{multiclass} task: how to represent these metrics with a single value (\emph{micro} and \emph{macro} averages)?
              \end{itemize}
    \end{enumerate}

    \begin{onlyenv}<3>
        \begin{tikzpicture}[overlay]
            \node[set, text width=34em, xshift=-2em, yshift=-3em] at (current page.center) {
                \begin{notebookin}
import numpy as np

means = np.array([[0.0, 0.0], [0.0, 3.0], [3.0, 0.0], [3.0, 3.0]])
y = np.repeat((0, 1, 2, 3), 250)  # 250 samples per class
X = means[y] + 0.6 * np.random.randn(4 * 250, 2)
print("X.shape =", X.shape)
print("y.shape =", y.shape)
                \end{notebookin}
                \begin{notebookout}
X.shape = (1000, 2)
y.shape = (1000,)
        \end{notebookout}
            };
        \end{tikzpicture}
    \end{onlyenv}

    \begin{onlyenv}<4>
        \begin{tikzpicture}[overlay]
            \node[set, xshift=-2em, yshift=-3em] at (current page.center) {
                \centering
                \includegraphics[width=0.55\textwidth]{02_slides/img/gaussian_data.png}
            };
        \end{tikzpicture}
    \end{onlyenv}

    \begin{onlyenv}<5>
        \begin{tikzpicture}[overlay]
            \node[set, text width=35em, xshift=-2em, yshift=1em] at (current page.center) {
                \begin{notebookin}
idx = np.random.permutation(X.shape[0])
X_train = X[idx[:800]]; y_train = y[idx[:800]]
X_test = X[idx[800:]]; y_test = y[idx[800:]]
                \end{notebookin}
            };
        \end{tikzpicture}
    \end{onlyenv}

    \begin{onlyenv}<6>
        \begin{tikzpicture}[overlay]
            \node[set, text width=35em, xshift=-2em, yshift=-3em] at (current page.center) {
                \begin{notebookin}
model = sm.MNLogit(y_train, sm.add_constant(X_train))
results = model.fit()
print(results.summary())
                \end{notebookin}

                \includegraphics[height=0.5\textheight]{02_slides/img/mnlogit.png}
            };
        \end{tikzpicture}
    \end{onlyenv}

    \begin{onlyenv}<7>
        \begin{tikzpicture}[overlay]
            \node[set, text width=35em, xshift=-2em, yshift=-3em] at (current page.center) {
                \begin{notebookin}
from sklearn.metrics import confusion_matrix

X_test = sm.add_constant(X_test)
y_pred = np.asarray(results.predict(X_test)).argmax(axis=1)
conf = confusion_matrix(y_test, y_pred)
                \end{notebookin}
            };
        \end{tikzpicture}
    \end{onlyenv}

    \begin{onlyenv}<8>
        \begin{tikzpicture}[overlay]
            \node[set, xshift=-2em, yshift=-3em] at (current page.center) {
                \centering
                \includegraphics[height=0.6\textheight]{02_slides/img/conf_mat.png}
            };
        \end{tikzpicture}
    \end{onlyenv}

    \begin{onlyenv}<9>
        \begin{tikzpicture}[overlay]
            \node[set, text width=37em, xshift=-2em, yshift=-3em] at (current page.center) {
                \begin{notebookin}
from sklearn.metrics import accuracy_score, precision_score, recall_score, f1_score

acc = accuracy_score(y_test, y_pred); print(f"Accuracy = {acc:.1%}")
for score in [precision_score, recall_score, f1_score]:
    for average in ["micro", "macro"]:
        metric = score(y_test, y_pred, average=average)
        print(f"{score.__name__.capitalize().removesuffix("_score")} ({average}) = {metric:.1%}")
                \end{notebookin}
                \begin{notebookout}
Accuracy = 98.5%
Precision (micro) = 98.5%
Precision (macro) = 98.4%
Recall (micro) = 98.5%
Recall (macro) = 98.6%
F1 (micro) = 98.5%
F1 (macro) = 98.5%
                \end{notebookout}
            };
        \end{tikzpicture}
    \end{onlyenv}
\end{frame}

\begin{frame}{Linear Models and Evaluation}{Preparing for next time\dots}
    You should now have mastered:
    \begin{itemize}
        \item How to prepare your data for ML (train/test splits, normalisation)
        \item What loss functions are and how to choose them
        \item How to define and use linear models for regression and classification
        \item How to regularise your models (and justify it)
    \end{itemize}
    \vfill \pause
    \begin{center}
        \huge
        Remember where you saved the data! Next time we shall start from there!
        \\[1em]
        Are there questions?
    \end{center}
    \vfill
\end{frame}
